%=====================================================================
% MundiPets — Parallel manuscript, version 1.0 (Deliverable 3 / 2A, week 13)
% Springer Nature journal template (sn-jnl.cls, v3.1 December 2024)
%
% Empirical component: APPROACH 2 of the course guide —
% automatic detection of ambiguity and smells in requirements.
%
% Closed sections: Introduction, Related Work, Methodology.
% Results, Discussion, Threats and Conclusions are declared placeholders
% and MUST NOT be written before Deliverable 4 (2B).
%
% Reference style: sn-mathphys-num (numbered). Switch the class option if
% the finally selected journal requires a different style.
%=====================================================================
\documentclass[pdflatex,sn-mathphys-num]{sn-jnl}

\usepackage{graphicx}%
\usepackage{multirow}%
\usepackage{amsmath,amssymb,amsfonts}%
\usepackage{amsthm}%
\usepackage{mathrsfs}%
\usepackage[title]{appendix}%
\usepackage{xcolor}%
\usepackage{textcomp}%
\usepackage{manyfoot}%
\usepackage{booktabs}%
\usepackage{algorithm}%
\usepackage{algorithmicx}%
\usepackage{algpseudocode}%
\usepackage{listings}%

% Visible marker for values the team must still fill in by hand.
\newcommand{\TODO}[1]{\textcolor{red}{\textbf{[TO COMPLETE: #1]}}}
% Marker for content deliberately deferred to Deliverable 4 (2B).
\newcommand{\DEFERRED}[1]{\textcolor{blue}{\textit{[Deferred to Deliverable 4 (2B): #1]}}}

\theoremstyle{thmstyleone}%
\newtheorem{theorem}{Theorem}%
\newtheorem{proposition}[theorem]{Proposition}%
\theoremstyle{thmstyletwo}%
\newtheorem{example}{Example}%
\newtheorem{remark}{Remark}%
\theoremstyle{thmstylethree}%
\newtheorem{definition}{Definition}%

\raggedbottom

\begin{document}

\title[Accuracy of a requirements ambiguity detector]{Accuracy of a
requirements ambiguity detector: case of a pet social network}

\author*[1]{\fnm{Edson Daniel} \sur{Fuertes Arraes}}\email{efuertesa2@uteq.edu.ec}

\author[1]{\fnm{Génesis Adriana} \sur{Gutiérrez Ortega}}\email{ggutierrezo@uteq.edu.ec}

\author[1]{\fnm{Andy Johel} \sur{Mendoza Párraga}}\email{amedozap9@uteq.edu.ec}

\author[1]{\fnm{Gary Alejandro} \sur{Morales Sánchez}}\email{gmoraless2@uteq.edu.ec}

\author[1]{\fnm{Jimmy Samuel} \sur{Nieves Sánchez}}\email{jnievess@uteq.edu.ec}

\author[1]{\fnm{Gleiston Cicerón} \sur{Guerrero Ulloa}}\email{gguerrero@uteq.edu.ec}

\affil*[1]{\orgdiv{Facultad de Ciencias de la Computación},
\orgname{Universidad Técnica Estatal de Quevedo},
\orgaddress{\street{Av. Quito km 1.5 vía a Santo Domingo},
\city{Quevedo}, \state{Los Ríos}, \country{Ecuador}}}

%%================================%%
%% Structured abstract             %%
%%================================%%
\abstract{\textbf{Context:} Ambiguity is one of the most persistent defects in
natural-language requirements specifications, and rule-based detectors are
frequently proposed as a low-cost mechanism to flag it before requirements reach
design. Evidence on how closely such detectors reproduce human expert judgement
remains scarce for Spanish-language specifications produced in Latin American
settings.

\textbf{Objective:} This study measures the extent to which a rule-based
ambiguity detector replicates the consensus judgement of independent human
experts over the complete requirements corpus of MundiPets, a real social
network for responsible pet adoption and breeding, and characterises the
patterns of disagreement between the two.

\textbf{Methods:} We conducted a paired quasi-experiment over 50 requirements
(25 functional, 16 non-functional and 9 design constraints) specified according
to an eight-attribute template. A purpose-built rule-based detector, implemented
in Python and frozen before data collection, classified each requirement as
ambiguous or non-ambiguous using four lexical and syntactic rules. In parallel,
three independent software engineers, blind to the detector output and to one
another, classified the same requirements in randomised order using a shared
rubric. The protocol was preregistered on the Open Science Framework before any
data were collected. Agreement was quantified with precision, recall, $F_1$,
Cohen's $\kappa$ and Fleiss' $\kappa$.

\textbf{Results:} \DEFERRED{quantitative results are reported in version 2.0 of
this manuscript, once the full statistical analysis described in Section~3.7 has
been executed through the versioned scripts.}

\textbf{Conclusion:} \DEFERRED{conclusions follow from the results.}}

\keywords{Requirements Engineering, Empirical Software Engineering, Natural
Language Processing, Interrater Agreement, Requirements Quality}

\maketitle

%=====================================================================
\section{Introduction}\label{sec:intro}

MundiPets is a social network for responsible pet adoption and breeding,
currently under specification for veterinary practices and pet owners in the
province of Los Ríos, Ecuador. The system lets owners publish verified pet
profiles, lets registered veterinary practices validate the associated medical
records, and supports adoption and breeding requests through a staged workflow
that includes compatibility assessment by an artificial-intelligence component.
Requirements engineering matters acutely in this domain because a defect in a
requirement about medical-record validation or about the visibility of an
owner's location does not merely degrade usability: it propagates into a system
that handles animal health data and personal data protected under Ecuador's
Organic Law on Personal Data Protection~\cite{lopdp2021}. Specifying such a
system with verifiable, unambiguous requirements is therefore a precondition for
lawful operation rather than a stylistic preference.

Ambiguity remains one of the defects most consistently reported in
natural-language requirements. The systematic mapping study of Montgomery et
al.~\cite{montgomery2022} shows that, although requirements quality has been
studied extensively, empirical work concentrates on a small number of quality
attributes and rarely validates automated detection against human judgement on
complete, real specifications. Atoum et al.~\cite{atoum2021} reach a convergent
conclusion in their systematic review of requirements quality assurance,
identifying the scarcity of validated automated support as an open challenge.
Ferrari et al.~\cite{ferrari2016} demonstrated that ambiguity in elicitation is
entangled with tacit knowledge, which suggests that purely lexical detection has
an intrinsic ceiling; yet the size of that ceiling has not been quantified for
specifications written in Spanish, nor for corpora produced by novice analysts,
the setting in which a large share of requirements documents are in fact
written. This is the gap the present study addresses.

This paper makes three contributions, each grounded in evidence collected by the
authors and released as an open dataset:
\begin{enumerate}
  \item a quantified comparison between a rule-based ambiguity detector and the
        blind consensus of three independent human experts over a complete
        requirements corpus of 50 items --- functional requirements,
        non-functional requirements and design constraints --- written in
        Spanish;
  \item a characterisation of the disagreement patterns between automated and
        human classification, reported per requirement and per rule, which
        exposes the specific linguistic phenomena that a rule-based approach
        systematically over- or under-flags;
  \item a replication package --- corpus, detector output, expert
        classifications, analysis scripts and preregistered protocol ---
        released under CC~BY~4.0 so that the study can be repeated on other
        corpora and other languages.
\end{enumerate}

The study answers two research questions:

\begin{itemize}
  \item[\textbf{RQ1}] With what precision, recall and $F_1$ does a rule-based
        automatic detector replicate the consensus judgement of human experts
        when classifying the requirements of MundiPets as ambiguous?
  \item[\textbf{RQ2}] Which disagreement patterns predominate between the
        automatic detector and expert judgement, and what do they reveal about
        the limits of a rule-based approach relative to human judgement in
        identifying ambiguity?
\end{itemize}

The remainder of the paper is organised as follows. Section~\ref{sec:related}
positions the study within the closest related work. Section~\ref{sec:method}
describes the experimental design, participants, instruments and analysis plan.
Section~\ref{sec:results} reports the results, Section~\ref{sec:discussion}
interprets them, Section~\ref{sec:threats} discusses threats to validity, and
Section~\ref{sec:conclusions} concludes.

%=====================================================================
\section{Related work}\label{sec:related}

This section is not a full systematic literature review; it follows the
abbreviated protocol of Kitchenham and Charters~\cite{kitchenham2007},
reporting a documented search string, explicit inclusion and exclusion criteria,
the number of studies retrieved and screened, and a comparative table of the
closest works.

\subsection{Search strategy}\label{subsec:search}

The search string combined three concept blocks:

\begin{quote}\small\ttfamily
("requirements engineering" OR "software requirements") AND
("ambiguity" OR "requirements smell" OR "requirements quality") AND
("automatic detection" OR "rule-based" OR "natural language processing")
\end{quote}

The string was executed on Scopus and Google Scholar on 2 August 2026,
restricted to the 2023--2026 period. Inclusion criteria: peer-reviewed
studies reporting empirical data on the quality, ambiguity or automated analysis
of natural-language requirements, published in requirements-engineering or
software-engineering venues. Exclusion criteria: editorials, short contributions
under four pages, and studies without empirical data. The search retrieved
689 records in Scopus and 3,670 in Google Scholar. After title and abstract
screening, 30 records were read in full text and 6 were retained. These were
combined with eight studies already identified in the project bibliography,
yielding the 14 entries compared in Table~\ref{tab:related}.

\subsection{Comparison with the closest studies}\label{subsec:comparison}

Table~\ref{tab:related} contrasts the retained studies with the present work.
The first eight rows correspond to studies already identified in the project
bibliography; the last six were retained from the search described above.

\begin{table}[h]
\caption{Comparison with the closest related studies}\label{tab:related}
\begin{tabular}{@{}p{2.1cm}p{0.7cm}p{2.2cm}p{2.4cm}p{4.0cm}@{}}
\toprule
Reference & Year & Study type & Population & Difference with our work \\
\midrule
Montgomery et al.~\cite{montgomery2022} & 2022 & Systematic mapping &
Primary studies on requirements quality & Maps the field but reports no primary
comparison between an automated detector and expert consensus on a single
complete corpus. \\
Atoum et al.~\cite{atoum2021} & 2021 & Systematic review & Literature on
requirements quality assurance and validation & Identifies the lack of validated
automation as an open challenge; does not measure detector--expert agreement
empirically. \\
Ferrari et al.~\cite{ferrari2016} & 2016 & Empirical study & Elicitation
interviews & Studies ambiguity as it arises in elicitation and its link to tacit
knowledge; our unit of analysis is the written requirement in a closed
specification. \\
Fischbach et al.~\cite{fischbach2023} & 2023 & Tool evaluation & Conditional
statements in natural-language requirements & Derives acceptance tests
automatically; shares the rule/NLP approach but targets test generation, not
ambiguity classification. \\
Gobov~\cite{gobov2023} & 2023 & Practitioner survey & Specification and
modelling techniques in practice & Characterises which techniques practitioners
use; provides no automated quality measurement. \\
Said et al.~\cite{said2026} & 2026 & Model evaluation & SRS documents &
Classifies requirements as functional or non-functional with a learned model;
our detector targets ambiguity and is rule-based and fully auditable. \\
Cheng et al.~\cite{cheng2026} & 2026 & Systematic review & Generative AI applied
to requirements engineering & Surveys generative approaches; we deliberately
avoid large language models to keep every classification traceable to a named
rule. \\
Vogelsang and Fischbach~\cite{vogelsang2024} & 2024 & Methodological guideline &
LLM use for NLP tasks in RE & Proposes usage guidelines rather than reporting an
agreement study. \\
Abdeahad et al.~\cite{abdeahad2026} & 2026 & Model evaluation &
7,061 requirements from the Fault-Prone SRS benchmark & Achieves high accuracy
with a deep model, but evaluates against dataset labels rather than against a
blind expert panel, and the model is not auditable rule by rule. \\
Fantechi et al.~\cite{fantechi2023} & 2023 & Tool comparison & Three example
requirements documents & Closest to our design: compares a rule-based tool with
an LLM on ambiguity. Uses the authors themselves as the reference judgement,
whereas our reference is a blind panel independent of the specification team. \\
Talha et al.~\cite{talha2025} & 2025 & Tool evaluation & 1,553 requirements
across 17 domains & Targets anaphoric and coordination ambiguity with a
transformer; our detector targets lexical and syntactic smells and is compared
against human consensus rather than a labelled benchmark. \\
Alemneh and Berhanu~\cite{alemneh2024} & 2024 & Systematic review & Literature
on requirements smells and detection techniques & Catalogues detection
techniques; reports no primary agreement measurement on a real corpus. \\
Intana et al.~\cite{intana2024} & 2024 & Tool evaluation & Thai-language SRS
documents & The closest non-English precedent; confirms that language-specific
evaluation is needed but does not cover Spanish nor report inter-rater
agreement. \\
Bashir et al.~\cite{bashir2025} & 2025 & Industrial study & Three railway
requirements datasets & Reports the over-flagging pattern (high recall, low
precision) also observed in rule-based tools; uses LLMs, which we deliberately
avoid for auditability. \\
\botrule
\end{tabular}
\end{table}

The niche of the present study is narrow and, to our knowledge, unoccupied: a
paired comparison between a frozen, fully documented rule-based detector and the
blind consensus of independent experts, over the \emph{complete} requirements
corpus of one real system, specified in Spanish, with the corpus, the detector
output and the individual expert classifications released as open data. Existing
work either maps the field without primary agreement data, or evaluates
detectors on English-language benchmark corpora that were not produced by the
team reporting the evaluation.

%=====================================================================
\section{Methodology}\label{sec:method}

The empirical component follows Approach~2 of the course framework ---
automatic detection of ambiguity and smells in requirements --- and is designed
as a paired quasi-experiment, reported following the guidelines of Jedlitschka
et al.~\cite{jedlitschka2008} and the experimental design principles of Wohlin
et al.~\cite{wohlin2012}. The protocol was preregistered on the Open Science
Framework~\cite{nosek2018} on 27 July 2026 at 10:49:41 (Ecuador time), before
the classification rubric was distributed to the expert panel and before the
detector was executed over the corpus.

\subsection{Research questions in PICOC format}\label{subsec:picoc}

\begin{table}[h]
\caption{Research questions structured with the PICOC framework}\label{tab:picoc}
\begin{tabular}{@{}p{2.0cm}p{9.5cm}@{}}
\toprule
Element & Description \\
\midrule
Population & The set of functional requirements (FR), non-functional
requirements (NFR) and design constraints (DC) specified for the MundiPets
system, written according to the eight-attribute template of the course and
current at the time the experiment was executed. \\
Intervention & Automatic classification of each requirement as ambiguous or
non-ambiguous by a rule-based detector that identifies vague quantifiers,
multiple conjunctions concealing more than one requirement, passive voice, and
absence of an explicit verification criterion. \\
Comparison & Manual and independent classification of the same requirement set
by three experts in requirements engineering, with no prior access to the
detector output (blind evaluation). \\
Outcome & Precision, recall and $F_1$ of the automatic detector against expert
consensus; inter-rater agreement via Cohen's $\kappa$ between pairs and Fleiss'
$\kappa$ across the three raters. \\
Context & Academic project of the Requirements Engineering course (Universidad
Técnica Estatal de Quevedo, 2026--2027 academic period) applied to the real
MundiPets system, with the specification written in Spanish. \\
\botrule
\end{tabular}
\end{table}

\subsection{Hypotheses}\label{subsec:hypotheses}

RQ1 is tested through the following pair of hypotheses:

\begin{description}
  \item[$H_0$] The classification produced by the automatic ambiguity detector
  is not associated with the consensus classification of the experts; the
  observed agreement between the two (Cohen's $\kappa$) is not significantly
  different from chance agreement ($\kappa = 0$).
  \item[$H_1$] The classification produced by the automatic detector is
  significantly associated with the expert consensus classification; the
  observed agreement is significantly greater than chance ($\kappa > 0$).
\end{description}

RQ2 is exploratory and is not subjected to hypothesis testing; it is answered
through qualitative analysis of the cases where detector and consensus diverge.

\subsection{Variables}\label{subsec:variables}

The \emph{independent variable} is the classification method applied to each
requirement, with two levels: (a) rule-based automatic detector and (b) consensus
of the three human experts. The \emph{dependent variable} is the resulting
binary classification of each requirement as ambiguous or non-ambiguous, from
which the agreement metrics are computed.

Four \emph{control variables} were fixed. First, requirement type (FR, NFR, DC)
was recorded to verify that detector performance does not vary systematically
with the artefact evaluated. Second, the detector configuration --- the four
rules listed in Section~\ref{subsec:design} --- remained invariant throughout
execution; no rule was adjusted after observing partial results. Third, the
presentation order of the requirements was randomised independently for each
rater to control fatigue and order effects. Fourth, raters were blind both to
the detector output and to one another's classifications until their own was
complete.

\subsection{Design and materials}\label{subsec:design}

The corpus comprises the complete set of 50 requirements of the MundiPets
specification: 25 functional requirements, 16 non-functional requirements and 9
design constraints. Each requirement was assigned an anonymous identifier
(REQ-01 to REQ-50), stripped of its type label, and presented to the raters in
randomised order. A key linking anonymous identifiers to the real requirement
identifiers and to the detector output was retained by the research team and not
shared with the raters.

The detector is a purpose-built Python program based on regular expressions. It
flags a requirement as ambiguous when at least one of four rules fires:
(i) presence of a vague quantifier; (ii) multiple conjunctions concealing more
than one requirement in a single statement; (iii) passive voice or absent
subject; and (iv) absence of an explicit verification criterion. No large
language model was used at any point of the classification pipeline, which keeps
every classification decision traceable to a named rule.

\subsection{Participants and recruitment}\label{subsec:participants}

The expert panel comprised three software engineers holding a professional
degree from the Universidad Técnica Estatal de Quevedo (two women and one man),
none of whom belonged to the MundiPets team or participated in writing the
requirements under evaluation. Inclusion criteria were: (a) a professional
degree in software engineering or a closely related field; (b) demonstrable
training or experience in requirements specification and analysis; and (c)
availability to complete the classification within the study schedule.

Each rater signed, on 28 July 2026, an informed consent form specific to this
role, drafted in compliance with Ecuador's Organic Law on Personal Data
Protection~\cite{lopdp2021}, explicitly authorising the anonymised use of their
classifications in peer-reviewed publications and their deposit in an open
dataset.

\subsection{Instruments}\label{subsec:instruments}

Two instruments were used. The \emph{expert classification rubric} is a workbook
containing an instruction sheet, a reference-criteria sheet for identifying
ambiguity, and an evaluation sheet where each rater records, per requirement,
the identifier, the full text, the classification, the type of smell detected, a
brief justification and a self-reported confidence level on a 1--5 scale. The
\emph{expert consent template} adapts the consent form used in the previous
elicitation round to the rater role, adding the clause on anonymised use in
peer-reviewed publications.

\subsection{Statistical analysis plan}\label{subsec:analysis}

Because the central variable is categorical, the primary analysis is based on
agreement metrics rather than mean comparison.

\bmhead{Primary analysis (RQ1, $H_0$/$H_1$)}
Precision, recall and $F_1$ of the detector are computed against expert
consensus, defined by majority vote across the three raters. The hypothesis is
tested through a significance test of Cohen's $\kappa$ (statistic
$z = \kappa / SE_\kappa$, against the standard normal distribution). The
magnitude of agreement is interpreted as the effect size using the Landis and
Koch scale: 0.00--0.20 slight; 0.21--0.40 fair; 0.41--0.60 moderate; 0.61--0.80
substantial; 0.81--1.00 almost perfect.

\bmhead{Inter-rater agreement}
Cohen's $\kappa$~\cite{cohen1960} is computed for each pair of raters and
Fleiss' $\kappa$~\cite{fleiss1971} across the three, interpreted on the same
scale. Low expert agreement would relativise any conclusion about detector
performance, since the reference consensus itself would be in doubt.

\bmhead{Normality test}
The Shapiro--Wilk test is applied to the continuous confidence variable
(1--5 scale) recorded in the rubric, as a precondition for the secondary
analysis.

\bmhead{Secondary analysis (RQ2, exploratory)}
To characterise disagreement, the confidence level reported by the raters is
compared between requirements where the detector agrees with consensus and those
where it diverges. If Shapiro--Wilk does not reject normality, an
independent-samples $t$ test is used with Cohen's $d$ as effect size; otherwise,
the Mann--Whitney $U$ test with Cliff's $\delta$. This is complemented by a
qualitative note on each disagreeing requirement.

\bmhead{Correction for multiple comparisons}
If agreement metrics are additionally reported disaggregated by requirement type
(FR, NFR, DC), the resulting $p$ values from those three related comparisons are
adjusted with the Benjamini--Hochberg procedure.

\bmhead{Power}
Following Molléri et al.~\cite{molleri2020} and Wohlin et al.~\cite{wohlin2012},
$\alpha = 0.05$ and $1-\beta = 0.80$ are fixed. Because the population size is
bounded by the real specification corpus rather than by an expandable
recruitment target, power is reported in reverse: instead of fixing a target
$N$, we determine the minimum detectable $\kappa$ that the design can identify
at that $\alpha$ and power, given the available $N = 50$.

\bmhead{Reproducibility}
All computations are executed through versioned analysis scripts; no figure or
table is produced manually.

%=====================================================================
\section{Results}\label{sec:results}

\DEFERRED{this section reports the analysed data without interpretation,
organised by research question, with every table and figure generated by the
versioned analysis scripts. It is written only after the full analysis plan of
Section~\ref{subsec:analysis} has been executed.}

%=====================================================================
\section{Discussion}\label{sec:discussion}

\DEFERRED{implications for research, implications for practice, and surprising
or counter-intuitive findings. Every claim must be traceable to a specific datum
in Section~\ref{sec:results} or to a cited reference.}

%=====================================================================
\section{Threats to validity}\label{sec:threats}

\DEFERRED{internal, external, construct and conclusion validity following Wohlin
et al.~\cite{wohlin2012}, with the mitigation applied and the residual
limitation acknowledged for each threat. The anticipated threats are already
documented in the preregistered protocol and are carried over here once the
results are known.}

%=====================================================================
\section{Conclusions and future work}\label{sec:conclusions}

\DEFERRED{one paragraph per research question, one consolidating the
contributions promised in Section~\ref{sec:intro}, and one with at least three
concrete lines of future work.}

%=====================================================================
\backmatter

\bmhead{Acknowledgements}

The authors thank the three software engineers who served as expert raters and
the veterinary practice in Quevedo, Ecuador, that granted access to its
operational context and which remains pseudonymised under the anonymisation
protocol of this study.

\section*{Declarations}

\begin{itemize}
\item \textbf{Funding.} This research received no specific grant from any
funding agency in the public, commercial or not-for-profit sectors.
\item \textbf{Competing interests.} The authors declare no competing interests.
\item \textbf{Ethics approval and consent to participate.} The study was
conducted under the ethics framework of the Requirements Engineering course at
the Universidad Técnica Estatal de Quevedo. All participants signed an informed
consent form drafted in compliance with Ecuador's Organic Law on Personal Data
Protection~\cite{lopdp2021}.
\item \textbf{Consent for publication.} All participants explicitly authorised
the use of their anonymised data in peer-reviewed publications and its deposit
in an open dataset.
\item \textbf{Data availability.} The preregistered protocol is available on the
Open Science Framework at https://osf.io/khyf2/overview The anonymised dataset --- labelled requirements corpus, detector output,
individual expert classifications, traceability matrix and analysis scripts ---
will be deposited on Zenodo under a CC~BY~4.0 licence with an assigned DOI upon
acceptance, following the FAIR principles~\cite{wilkinson2016} and the
software citation principles of Smith et al.~\cite{smith2016}.
\item \textbf{Materials availability.} Not applicable.
\item \textbf{Code availability.} The detector and the analysis scripts are
available in the project repository at
https://github.com/andymendozap03/MundiPets-PFC-IngReq
\item \textbf{Author contribution.} E.D.F.A.: conceptualization,
investigation, writing --- original draft. G.A.G.O.: software (design and
implementation of the rule-based detector). A.J.M.P.: methodology,
visualization (system modelling and diagrams). G.A.M.S.: investigation, data
curation (interviews and transcriptions). J.S.N.S.: formal analysis, data
curation, project administration (protocol preregistration and statistical
analysis). G.C.G.U.: supervision, validation. All authors reviewed and approved
the final manuscript.
\end{itemize}

\bibliography{referencias}

\end{document}
